%% 6-7 ゆれながら止まる：減衰振動（第14回：組 V6-4。ページ制限なし。偶数ページ）
%% ほかの巻が指している形：5-16「ゆれながら止まる動き（減衰振動、6-7・6-12）」、5-18「ゆれながら止まる（e^{at} と sin のかけ算）：
%%   計器の針、車のサスペンション（6-7・6-12）」。
%% 記号：減衰係数 c〔N·s/m〕（6-5 の2乗の抵抗の c とは別の量）、減衰比 ζ=c/(2√(mk))、へる速さ α=c/(2m)=ζω_n、
%%   対数減衰率 δ（となりの山の比の自然対数）。ζ の理由（特性方程式）は 6-12。
%% 同じ数：ばね 0.5 kg・8 N/m、c=0.4・4・10。RLC は 10 mH・1 μF、R=20 Ω・200 Ω・1 kΩ。表計算は 6-2 の型（速度を先に進めるオイラー法）
\makeatletter
\newcommand\kvihhead[1]{\par\addvspace{6pt}\Needspace{6\baselineskip}\noindent{\sffamily\bfseries\color{\bkTT}#1}\par\nobreak\vspace{2pt}}
\newcommand\kvihbox[1]{\par\addvspace{3pt}\noindent\textbf{#1}\par\noindent\ignorespaces}
\newcommand\kviha{\par\vspace{2pt}\noindent\hangindent=2.6\zw\hangafter=0\relax}
\newcommand\kvihb{\par\vspace{2pt}\noindent\hangindent=4.6\zw\hangafter=0\relax}
\newcommand\kvihframe[1]{\par\addvspace{4pt}{\centering\tikz\node[draw=\bkTT, line width=0.9pt, fill=bkTheme!8, rounded corners=1.5mm, inner xsep=6mm, inner ysep=2.5mm, align=center]{#1};\par}\addvspace{4pt}\noindent\ignorespaces}
% グラフの目もりの数字を、グラフの下のはしの外に置く  \kvihxtl{x}{ymin}{数字}
\newcommand\kvihxtl[3]{\node[font=\sffamily\footnotesize, text=bkInk, anchor=north, inner sep=1pt] at (axis cs:#1,#2) {#3};}
% ばね（横）  \kvihspringh{左の x}{右の x}{y}
\newcommand\kvihspringh[3]{%
  \draw[bkInk!80, line width=0.7pt] (#1,#3) -- ({#1+2},#3);
  \draw[bkInk!80, line width=0.7pt, decorate, decoration={zigzag, segment length=2mm, amplitude=1.7mm}] ({#1+2},#3) -- ({#2-2},#3);
  \draw[bkInk!80, line width=0.7pt] ({#2-2},#3) -- (#2,#3);}
% ダンパー（横）  \kvihdamper{左の x}{右の x}{y}：左に油の入った筒、右からピストンの棒
\newcommand\kvihdamper[3]{%
  \draw[bkInk!80, line width=0.7pt] (#1,#3) -- ({#1+2},#3);
  \filldraw[fill=bkYellow!35, draw=bkInk!80, line width=0.7pt] ({#1+2},{#3-2.6}) rectangle ({(#1+#2)/2+3},{#3+2.6});
  \draw[bkInk!80, line width=0.7pt] ({(#1+#2)/2+3},{#3-2.6}) -- ({(#1+#2)/2+3},{#3+2.6});
  \fill[bkInk!80] ({(#1+#2)/2-1.2},{#3-2.2}) rectangle ({(#1+#2)/2},{#3+2.2});
  \draw[bkInk!80, line width=0.9pt] ({(#1+#2)/2},#3) -- (#2,#3);}
% ダンパー（たて）  \kvihdamperv{x}{上の y}{下の y}：下に筒、上からピストンの棒
\newcommand\kvihdamperv[3]{%
  \draw[bkInk!80, line width=0.7pt] (#1,#3) -- (#1,{#3+2});
  \filldraw[fill=bkYellow!35, draw=bkInk!80, line width=0.7pt] ({#1-2.6},{#3+2}) rectangle ({#1+2.6},{(#2+#3)/2+3});
  \draw[bkInk!80, line width=0.7pt] ({#1-2.6},{(#2+#3)/2+3}) -- ({#1+2.6},{(#2+#3)/2+3});
  \fill[bkInk!80] ({#1-2.2},{(#2+#3)/2-1.2}) rectangle ({#1+2.2},{(#2+#3)/2});
  \draw[bkInk!80, line width=0.9pt] (#1,{(#2+#3)/2}) -- (#1,#2);}
\newcommand\kvihspringv[3]{%
  \draw[bkInk!80, line width=0.7pt] (#1,#2) -- (#1,{#2-2});
  \draw[bkInk!80, line width=0.7pt, decorate, decoration={zigzag, segment length=2mm, amplitude=1.7mm}] (#1,{#2-2}) -- (#1,{#3+2});
  \draw[bkInk!80, line width=0.7pt] (#1,{#3+2}) -- (#1,#3);}
\newcommand\kvihwall[3]{%
  \fill[pattern=north east lines, pattern color=bkGray] ({#1-2.6},#2) rectangle (#1,#3);
  \draw[bkInk, line width=1pt] (#1,#2) -- (#1,#3);}
\newcommand\kvihfloor[3]{%
  \fill[pattern=north east lines, pattern color=bkGray] (#1,{#3-2.4}) rectangle (#2,#3);
  \draw[bkInk, line width=1pt] (#1,#3) -- (#2,#3);}
\newcommand\kvihcart[2]{\begin{scope}[shift={#1}]
  \filldraw[fill=#2!25, draw=#2!80!black, line width=0.8pt, rounded corners=1mm] (-7,3) rectangle (7,15);
  \filldraw[fill=bkInk!70, draw=bkInk, line width=0.5pt] (-4,1.5) circle (1.5) (4,1.5) circle (1.5);
\end{scope}}
% 回路の部品（たて向き）
\newcommand\kvihcap[3]{%
  \draw[bkInk, line width=0.8pt] (#1,#2) -- (#1,{(#2+#3)/2-1.1}) (#1,{(#2+#3)/2+1.1}) -- (#1,#3);
  \draw[bkInk, line width=1.4pt] ({#1-4},{(#2+#3)/2-1.1}) -- ({#1+4},{(#2+#3)/2-1.1}) ({#1-4},{(#2+#3)/2+1.1}) -- ({#1+4},{(#2+#3)/2+1.1});}
\newcommand\kvihcoil[3]{%
  \draw[bkInk, line width=0.8pt] (#1,#2) -- (#1,{(#2+#3)/2-6}) (#1,{(#2+#3)/2+6}) -- (#1,#3);
  \draw[bkInk, line width=0.8pt, decorate, decoration={coil, aspect=0.55, segment length=2.4mm, amplitude=2.4mm}] (#1,{(#2+#3)/2-6}) -- (#1,{(#2+#3)/2+6});}
\newcommand\kvihres[3]{%
  \draw[bkInk, line width=0.8pt] (#1,#2) -- (#1,{(#2+#3)/2-5}) (#1,{(#2+#3)/2+5}) -- (#1,#3);
  \draw[bkInk, line width=0.8pt, fill=white] ({#1-1.8},{(#2+#3)/2-5}) rectangle ({#1+1.8},{(#2+#3)/2+5});}
\makeatother

\bkchapter{6}{7}{ゆれながら止まる：\\減衰振動}{\bklevelrange{3}{4}}{ゆれながら止まる：減衰振動}
\bkmichishirube{3-11,5-14,5-15,6-2,6-5,6-6}{6-12,6-13,6-15}

\begin{goal}
\begin{itemize}[label={\color{\bkTT}●}, leftmargin=1.6\zw, itemsep=0pt, topsep=1pt]
\item \bklevel{3}\enspace ばね・質量・ダンパーの式 $m\ddot x+c\dot x+kx=0$ を、力の図から立てられる。
\item \bklevel{3}\enspace 表計算で動かし、となりの山の比から、ゆれのへり方（対数減衰率）を読みとれる。
\item \bklevel{4}\enspace 減衰比 $\zeta$ を計算して、「ゆれる・ちょうど・ゆっくり」のどれになるかを答えられる（RLC回路でも）。
\end{itemize}
\end{goal}

\begin{hatena}
\noindent\begin{minipage}[c]{0.56\linewidth}
\raggedright
理科室の電流計（アナログのメーター）にスイッチを入れると、針は目もりまでさっと動き、ほとんど行きすぎずに、すっと止まります。もし針がばねだけで動いていたら、目もりのまわりをいつまでもゆれ続けて、読めません。

針には、ゆれを止めるブレーキがついています。ブレーキは、どのくらいの強さにすればよいのでしょう。弱すぎても、強すぎても困りそうです。
\end{minipage}\hfill
\begin{minipage}[c]{0.41\linewidth}
\centering
\begin{tikzpicture}[x=1mm, y=1mm]
  \filldraw[fill=white, draw=bkInk, line width=0.8pt, rounded corners=1.5mm] (-26,-6) rectangle (26,26);
  \draw[bkInk, line width=0.7pt] (130:20) arc (130:50:20);
  \foreach \a in {130,120,...,50} {\draw[bkInk, line width=0.6pt] (\a:20) -- (\a:17.5);}
  \foreach \a/\l in {130/0, 90/5, 50/10} {\node[font=\sffamily\footnotesize, text=bkInk] at (\a:23.2) {\l};}
  \draw[bkRed!70, line width=0.5pt, dashed] (0,0) -- (78:16.5);
  \draw[bkRed!70, line width=0.5pt, dashed] (0,0) -- (102:16.5);
  \draw[bkInk, line width=1.1pt] (0,0) -- (90:16.5);
  \fill[bkInk] (0,0) circle (1.4);
  \node[font=\sffamily\footnotesize, text=bkInk] at (0,-3.4) {A};
  \node[font=\sffamily\footnotesize, anchor=north, align=center] at (0,-7.5) {ゆれ続けたら\\目もりが読めない};
\end{tikzpicture}
\end{minipage}
\end{hatena}

\noindent 6-6では、ゆれは止まらないと考えました。でも本物のゆれは、空気の抵抗や摩擦で、いつか止まります。車や建物や計器では、わざと「ゆれを止める部品」をつけます。この章では、その部品（ダンパー）を式に入れて、ゆれがどのように止まるのかを、表計算で調べます。

%% ---------------- やってみよう・思い出そう ----------------
\begin{yattemiyou}[ゆれの山の高さをはかろう]
\kvihbox{用意するもの}
6-6のやってみようのばねとおもり（\qty{100}{g}）、水を入れた大きなペットボトル（上を切ったもの）、ものさし、スマートフォン（スロー撮影）
\begin{enumerate}[label={\bkcircled[bkGreen]{\arabic*}}, leftmargin=2.4\zw, itemsep=1pt, topsep=3pt]
\item ものさしをおもりの横に立て、つり合いの位置から\qty{5}{cm}引いてはなし、スロー撮影しよう。
\item 動画をコマ送りして、上に来たときの「つり合いの位置からの高さ」（山の高さ）を、1回めから順に読もう。
\item おもりを水の中に入れて、①②をくり返そう。
\item となりの山の高さの比（あとの山÷前の山）を計算しよう。
\end{enumerate}
\begin{bkcenter}
\begin{bktable}{colspec={Q[l,m,wd=28mm] *{4}{Q[c,m,wd=17mm]}}, rowsep=1pt}
山 & 1回め & 2回め & 3回め & 4回め \\
空気の中（cm） & 5.0 & 4.6 & 4.2 & 3.9 \\
水の中（cm） & 5.0 & 2.0 & 0.8 & 0.3 \\
\end{bktable}
\bktabcaption{ソラさんがはかった結果}
\end{bkcenter}
\end{yattemiyou}

\noindent となりの山の比（あとの山÷前の山）を計算します。
\begin{itemize}[leftmargin=1.2\zw, itemsep=1pt, topsep=2pt, label={\color{\bkTT}●}]
\item 空気の中：$4.6\div5.0=0.92$、\quad $4.2\div4.6\approx0.91$、\quad $3.9\div4.2\approx0.93$
\item 水の中：$2.0\div5.0=0.4$、\quad $0.8\div2.0=0.4$、\quad $0.3\div0.8\approx0.38$
\end{itemize}

どちらも、\textbf{となりの山の比が、ほぼ同じ}になりました。山の高さは「同じ数ずつ」ではなく、「同じ割合ずつ」へっています。3-11で学んだ「倍々でへる」形です。水の中のほうが、1回でへる割合が大きく、早く止まります。

\begin{omoidasou}
\bkrecall{3-11}{同じ割合ずつへる量は、指数関数で表せる（倍々でへる）。}
\bkrecall{5-14}{$e^{-t/\tau}$ は、同じ時間ごとに同じ割合でへる。$e^x=y$ なら $x=\ln y$。}
\bkrecall{6-2}{2階の式は、位置 $x$ と速度 $v$ の2つの列で進める（速度を先に進めるオイラー法）。}
\bkrecall{6-5}{空気の抵抗は、動きと反対向き。速さに比例する抵抗は $f=bv$。\\2乗に比例する抵抗の係数を $c$ と書いた（$f=cv^2$）。}
\bkrecall{6-6}{ばねとおもり：$m\ddot x=-kx$、$\omega_{\mathrm n}=\sqrt{\dfrac{k}{m}}$。\qty{0.5}{kg}・\qty{8}{N/m}で $\omega_{\mathrm n}=4$、周期\qty{1.57}{s}。\\エネルギー $\dfrac{1}{2}mv^2+\dfrac{1}{2}kx^2$。LC回路は $L\ddot q=-\dfrac{1}{C}q$。}
\end{omoidasou}

%% ---------------- 数学の道具 ----------------
\bksection{doguu}{数学の道具}

\bksubsection{1}{ダンパー：速さに比例するブレーキ}

\label{bk:6-7:damper}ゆれを止めるために使う部品を、\textbf{ダンパー}といいます。車のタイヤの近くについている「ショックアブソーバー」や、ドアがゆっくり閉まるようにする「ドアクローザー」もダンパーです。\bkidx{だんぱー}{ダンパー}

\begin{bkfigure}
\begin{tikzpicture}[x=1mm, y=1mm]
  % 筒とピストン（大きく）
  \filldraw[fill=bkYellow!30, draw=bkInk, line width=0.9pt] (0,0) rectangle (60,20);
  \draw[bkInk, line width=0.9pt] (60,0) -- (60,20);
  \fill[bkInk!75] (34,1) rectangle (38,19);
  \fill[bkYellow!30] (34,5) rectangle (38,7) (34,13) rectangle (38,15);
  \draw[bkInk!80, line width=2pt] (38,10) -- (84,10);
  \draw[bk arrow, bkBlue, line width=1.2pt] (70,15) -- (82,15);
  \node[font=\sffamily\footnotesize, anchor=south, text=bkBlue] at (76,15.8) {棒を引く};
  % 油の流れ
  \draw[bk arrow, bkOrange!85!black, line width=0.8pt] (42,6) -- (30,6);
  \draw[bk arrow, bkOrange!85!black, line width=0.8pt] (42,14) -- (30,14);
  \node[font=\sffamily\footnotesize, anchor=north, align=center] at (17,-1.5) {油};
  \node[font=\sffamily\footnotesize, anchor=north, align=center] at (36,-1.5) {小さな穴の\\あいた板};
  \draw[bkGray, line width=0.4pt] (36,-1.8) -- (36,1);
  \node[font=\sffamily\footnotesize, anchor=west, align=left, text width=40mm] at (88,10) {速く引くほど、油は小さな穴を通りにくく、強くじゃまをする};
\end{tikzpicture}
\bkfigcaption{ダンパーのしくみ。筒の中に油が入っていて、小さな穴のあいた板（ピストン）が動く。}
\end{bkfigure}

筒の中の油が、ピストンの小さな穴を通りぬけないと、棒は動けません。ゆっくり引けば油はゆっくり流れればよいので、力はほとんどいりません。速く引くと、油を急いで流さなければならないので、強い力が必要です。ダンパーの力の大きさは、ほぼ速さに比例します。

\kvihhead{力の向きは、速度と反対}
ダンパーは、動いている向きと反対向きにじゃまをします。右へ動いているときは左へ、左へ動いているときは右へ引きもどします。速度を $v=\dot x$（右向きが＋）とすると、
\kvihframe{ダンパーの力\qquad $-c\dot x$}
比例定数 $c$ を\textbf{減衰係数}といいます。単位は、力〔N〕÷速さ〔m/s〕なので〔N·s/m〕です。たとえば $c=\qty{0.4}{N.s/m}$ のダンパーは、\qty{0.2}{m/s}で右へ動いているとき、$-0.4\times0.2=\qty{-0.08}{N}$（左向きに\qty{0.08}{N}）の力を出します。

\begin{seikaku}
6-5では、空気の2乗に比例する抵抗の係数を $c$（単位〔N·s$^2$/m$^2$〕）と書きました。この章からは、ダンパーの減衰係数を $c$（単位〔N·s/m〕）と書きます。同じ文字ですが、単位のちがう別の量です。機械の本では、ダンパーの係数を $c$ と書くのがふつうなので、この本もそれにそろえます。ダンパーの力は、6-5の「速さに比例する抵抗 $bv$」と同じ形です。
\end{seikaku}

\bksubsection{2}{ばね・質量・ダンパーの式を立てる}

\label{bk:6-7:shiki}6-6の台車（\qty{0.5}{kg}、ばね\qty{8}{N/m}）に、ダンパーをならべてつけます。

\begin{bkfigure}
\begin{tikzpicture}[x=1mm, y=1mm]
  \kvihfloor{0}{110}{0}
  \kvihwall{4}{0}{24}
  \kvihspringh{4}{50}{14}
  \kvihdamper{4}{50}{5.5}
  \kvihcart{(57,0)}{bkBlue}
  \node[font=\sffamily\footnotesize, anchor=south] at (27,17.5) {ばね $k$};
  \node[font=\sffamily\footnotesize, anchor=north] at (27,1.8) {};
  \node[font=\sffamily\footnotesize, anchor=south] at (57,15.5) {質量 $m$};
  % 力の矢印（右へ動いているとき、x>0）
  \draw[bk arrow, bkRed, line width=1.1pt] (64,11) -- (78,11);
  \node[font=\sffamily\footnotesize, anchor=west, text=bkRed] at (79,11) {速度 $v$（右向き）};
  \draw[bkGray, dashed, line width=0.5pt] (57,-3) -- (57,26);
  \node[font=\sffamily\footnotesize, anchor=south] at (57,26.5) {$x=0$};
\end{tikzpicture}
\bkfigcaption{ばね（上）とダンパー（下）が、ならんで台車をかべにつないでいる。}
\end{bkfigure}

\kvihhead{台車にはたらく力}
台車が位置 $x$ にいて、速度 $\dot x$ で動いているとき、左右の向きの力は2つです。
\begin{itemize}[leftmargin=1.2\zw, itemsep=2pt, topsep=2pt, label={\color{\bkTT}●}]
\item ばねの力 $-kx$（ずれと反対向き。6-6）
\item ダンパーの力 $-c\dot x$（速度と反対向き。①）
\end{itemize}
$F=ma$（5-5）に入れると、
\begin{bkhenkei}
m\ddot x=-kx-c\dot x & 質量×加速度＝力の合計 \\
m\ddot x+c\dot x+kx=0 & 右の2つを左へ移す \\
\end{bkhenkei}
\kvihframe{$m\ddot x+c\dot x+kx=0$}
これが、\textbf{ばね・質量・ダンパー}の式です。左から「質量×加速度」「ダンパー」「ばね」の3つが、1つずつ入っています。車のサスペンションも、建物のゆれも、計器の針も、まずこの形で考えます。

\kvihhead{2つの力の向きをくらべる}
ばねの力は「位置」と反対、ダンパーの力は「速度」と反対です。ゆれる台車の4つの場面で、向きを見ます（$c=0.4$、振幅\qty{0.1}{m}のとき）。
\begin{bkcenter}
\begin{bktable}{colspec={Q[l,m,wd=34mm] Q[c,m,wd=20mm] Q[c,m,wd=20mm] Q[c,m,wd=22mm] Q[c,m,wd=22mm]}, rowsep=1pt}
場面 & 位置 $x$ & 速度 $\dot x$ & ばね $-kx$ & ダンパー $-c\dot x$ \\
右のはし & ＋ & 0 & 左向き & 0 \\
まん中を左へ & 0 & － & 0 & 右向き \\
左のはし & － & 0 & 右向き & 0 \\
まん中を右へ & 0 & ＋ & 0 & 左向き \\
\end{bktable}
\bktabcaption{ばねの力は、はしで最大。ダンパーの力は、まん中（速さが最大）で最大。}
\end{bkcenter}
ダンパーは、はしでは何もしません（止まっているから）。まん中を走りぬけるときに、いちばん強くブレーキをかけます。

\kvihhead{答えの候補の表には、ない形}
5-18の答えの候補の表には、1回微分と2回微分が両方入った式はありません。ためしに $x=\sin4t$ を入れてみます（$m=0.5$、$c=0.4$、$k=8$）。
\begin{bkhenkei}
0.5\ddot x+8x=-8\sin4t+8\sin4t=0 & 質量とばねの項（6-6と同じ）\\
0.4\dot x=0.4\times4\cos4t=1.6\cos4t & ダンパーの項 \\
\end{bkhenkei}
合計は $1.6\cos4t$ になり、0になりません。ダンパーの項が、$\sin$ を微分した $\cos$ を生んでしまい、ほかの項と打ち消しあえないのです。$\sin$ だけでも、$e^{at}$ だけでも、うまくいきません。そこで、まず表計算で動きを見ます。式で解く方法は、6-12で学びます。

\bksubsection{3}{表計算で動かしてみる}

\label{bk:6-7:hyou}6-2・6-6と同じ型で計算します。変わるのは、加速度の式だけです。
\begin{enumerate}[label={\bkcircled[\bkTF]{\arabic*}}, leftmargin=2.4\zw, itemsep=1pt, topsep=2pt]
\item 今の位置と今の速度から、加速度を求める：$a=\dfrac{-kx-cv}{m}$（ここが6-6とちがう）
\item 新しい速度：$v+a\times\Delta t$
\item 新しい位置：$x+(\text{新しい速度})\times\Delta t$
\end{enumerate}
①の式は、②の運動方程式 $m\ddot x=-kx-c\dot x$ の両辺を $m$ でわったものです。
$c=0.4$、$\Delta t=0.01$ として、D2に \texttt{=(-8*B2-0.4*C2)/0.5}、C3に \texttt{=C2+D2*\$E\$2}、B3に \texttt{=B2+C3*\$E\$2} と書いて、下へコピーします。
\begin{bkcenter}
\bksheet{5}{19}{%
1 & $t$〔s〕 & $x$〔m〕 & $v$〔m/s〕 & $a$〔m/s$^2$〕 & $\Delta t$ \\
2 & 0 & 0.10000 & 0 & $-1.60000$ & 0.01 \\
3 & 0.01 & 0.09984 & $-0.01600$ & $-1.58464$ & \\
4 & 0.02 & 0.09952 & $-0.03185$ & $-1.56687$ & \\
81 & 0.79 & $-0.07286$ & 0.00265 & 1.16369 & \\
159 & 1.57 & 0.05312 & 0.00468 & $-0.85374$ & \\}
\bktabcaption{4行めのあとは、とちゅうの行を省いた。81行め（0.79秒）のあたりで速度の符号が変わる（左のはし）。}
\end{bkcenter}
2行めの加速度は、速度が0なので、6-6と同じ $-16\times0.1=-1.6$ です。3行めからは、ダンパーの $-0.4\times v$ が加わります。たとえば3行めの加速度は
\[ \frac{-8\times0.09984-0.4\times(-0.016)}{0.5}=\frac{-0.79872+0.0064}{0.5}=-1.58464 \]
で、6-6（$-1.59744$）より少し小さくなります。左へ動いているので、ダンパーが右向きに引いて、加速をじゃましているのです。

\kvihhead{3つのダンパーで、くらべる}
$c$ を0.4、4、10にかえて、同じ計算をします（どれも、はじめ $x=0.1$、$v=0$）。
\begin{bkfigure}
\begin{tikzpicture}
\begin{axis}[bkplot, width=124mm, height=54mm, xmin=0, xmax=6.3, ymin=-0.16, ymax=0.118, xtick={1,2,3,4,5,6}, xticklabels={}, ytick={-0.05,0.05,0.1}, yticklabels={$-0.05$,0.05,0.1}, xlabel={$t$（s）}, ylabel={$x$（m）}]
  \addplot[bkcurve, domain=0:6.1, samples=220] {0.1*exp(-0.4*x)*(cos(deg(3.97995*x))+0.100504*sin(deg(3.97995*x)))};
  \addplot[bkcurve2, domain=0:6.1, samples=120] {0.1*(1+4*x)*exp(-4*x)};
  \addplot[bkPurple, line width=1.7pt, domain=0:6.1, samples=120] {0.104564*exp(-0.834849*x)-0.004564*exp(-19.165151*x)};
  \kvihxtl{1}{-0.16}{1}\kvihxtl{2}{-0.16}{2}\kvihxtl{3}{-0.16}{3}\kvihxtl{4}{-0.16}{4}\kvihxtl{5}{-0.16}{5}\kvihxtl{6}{-0.16}{6}
  \node[font=\sffamily\footnotesize, text=bkBlue, anchor=west, inner sep=1pt, fill=white] at (axis cs:2.7,-0.06) {$c=0.4$：ゆれながら止まる};
  \node[font=\sffamily\footnotesize, text=bkOrange!80!black, anchor=west, inner sep=1pt, fill=white] at (axis cs:2.7,-0.092) {$c=4$：ゆれずに、早くもどる};
  \node[font=\sffamily\footnotesize, text=bkPurple, anchor=west, inner sep=1pt, fill=white] at (axis cs:2.7,-0.124) {$c=10$：ゆっくりもどる};
\end{axis}
\end{tikzpicture}
\bkfigcaption{ダンパーの強さで、止まり方が3つに分かれる（青 $c=0.4$、オレンジ $c=4$、紫 $c=10$）。}
\label{bk:6-7:fig-three}
\end{bkfigure}
\begin{itemize}[leftmargin=1.2\zw, itemsep=2pt, topsep=2pt, label={\color{\bkTT}●}]
\item \textbf{$c=0.4$（弱いダンパー）}\quad ゆれながら、山がだんだん低くなって止まる。
\item \textbf{$c=4$}\quad 0をこえずに、すっともどる。3本の中で、いちばん早く0に近づく。
\item \textbf{$c=10$（強いダンパー）}\quad 0をこえないが、ゆっくりとしかもどらない。2秒たっても、まだ約\qty{2}{cm}残っている。
\end{itemize}
強いダンパーほど早く止まる、とは限りません。強すぎるダンパーは、もどる動きまでじゃましてしまうからです。

\bksubsection{4}{山の高さは、同じ割合でへる}

\label{bk:6-7:yama}$c=0.4$ のグラフの山（右のはし）を、表計算から読みとります。
\begin{bkcenter}
\begin{bktable}{colspec={Q[l,m,wd=36mm] *{5}{Q[c,m,wd=15mm]}}, rowsep=1pt}
山 & 1回め & 2回め & 3回め & 4回め & 5回め \\
時刻（s） & 0 & 1.579 & 3.157 & 4.736 & 6.314 \\
山の高さ（m） & 0.1 & 0.0532 & 0.0283 & 0.0150 & 0.0080 \\
前の山との比 & & 0.532 & 0.532 & 0.532 & 0.532 \\
\end{bktable}
\bktabcaption{山と山の時間は、どれも1.579秒。となりの山の比は、どれも0.532。}
\end{bkcenter}
やってみようと同じように、\textbf{となりの山の比が一定}です。山は、1.579秒ごとに0.532倍になる「倍々でへる」形です。

\kvihhead{つつみの曲線}
「同じ時間ごとに同じ割合でへる」量は、指数関数 $e^{-\alpha t}$ で書けます（5-14）。山と山の時間 $T_{\mathrm d}=1.579$ で0.532倍になるので、
\begin{bkhenkei}
e^{-\alpha\times1.579}=0.532 & 1.579秒で0.532倍 \\
-\alpha\times1.579=\ln0.532=-0.631 & 両辺の自然対数をとる（5-14）\\
\alpha=\frac{0.631}{1.579}\approx0.400 & 両辺を $-1.579$ でわる \\
\end{bkhenkei}
山の頂上を結ぶと、$0.1e^{-0.4t}$ の曲線になります。$T_{\mathrm d}$ の d は、damped（減衰した）の頭文字です（6-12の $\omega_{\mathrm d}$ と同じ）。この曲線を、ゆれの\textbf{つつみ}といいます。ゆれは、上下のつつみ $\pm0.1e^{-0.4t}$ のあいだにおさまりながら、小さくなっていきます。ゆれそのものは、つつみ $0.1e^{-0.4t}$ と $\cos$ の波のかけ算に近い形です（5-18⑥の「ゆれながら止まる：$e^{at}$ と $\sin$ のかけ算」の行。くわしくは6-12）。

\begin{bkfigure}
\begin{tikzpicture}
\begin{axis}[bkplot, width=124mm, height=50mm, xmin=0, xmax=6.5, ymin=-0.11, ymax=0.125, xtick={1.579,3.157,4.736,6.314}, xticklabels={}, ytick={-0.1,0.1}, yticklabels={$-0.1$,0.1}, xlabel={$t$（s）}, ylabel={$x$（m）}]
  \addplot[bkthin, dashed, domain=0:6.3, samples=60] {0.1*exp(-0.4*x)};
  \addplot[bkthin, dashed, domain=0:6.3, samples=60] {-0.1*exp(-0.4*x)};
  \addplot[bkcurve, domain=0:6.3, samples=220] {0.1*exp(-0.4*x)*(cos(deg(3.97995*x))+0.100504*sin(deg(3.97995*x)))};
  \addplot[only marks, mark=*, mark size=1.6pt, bkOrange] coordinates {(0,0.1) (1.579,0.0532) (3.157,0.0283) (4.736,0.0150) (6.314,0.0080)};
  \kvihxtl{1.579}{-0.11}{1.58}\kvihxtl{3.157}{-0.11}{3.16}\kvihxtl{4.736}{-0.11}{4.74}\kvihxtl{6.314}{-0.11}{6.31}
  \node[font=\sffamily\footnotesize, text=bkGray, anchor=south west, fill=white, inner sep=1pt] at (axis cs:2.0,0.052) {つつみ $0.1e^{-0.4t}$};
\end{axis}
\end{tikzpicture}
\bkfigcaption{山（オレンジの点）は、つつみの曲線（点線）の上にのる。}
\end{bkfigure}

\kvihhead{へる速さ $\alpha$ は、$c$ と $m$ で決まる}
$\alpha=0.4$ はどこから来たのでしょう。$c$ だけを0.2、1、2にかえて同じことをすると、次のようになります（$m=0.5$）。
\begin{bkcenter}
\begin{bktable}{colspec={Q[l,m,wd=40mm] *{4}{Q[c,m,wd=17mm]}}, rowsep=1pt}
減衰係数 $c$ & 0.2 & 0.4 & 1 & 2 \\
つつみから読んだ $\alpha$ & 0.2 & 0.4 & 1.0 & 2.0 \\
$\dfrac{c}{2m}$ & 0.2 & 0.4 & 1.0 & 2.0 \\
\end{bktable}
\bktabcaption{$\alpha$ はいつも $\dfrac{c}{2m}$ と同じ。}
\end{bkcenter}
ダンパーが強いほど速くへり、台車が重いほどゆっくりへる、というきまりです。
\kvihframe{つつみ\quad $e^{-\alpha t}$\qquad $\alpha=\dfrac{c}{2m}$}
なぜ2でわるのかは、このあとの「もっと知りたい」でエネルギーから説明します。

\kvihhead{山と山の時間は、少し長くなる}
山と山の時間 $T_{\mathrm d}=1.579$ 秒は、ダンパーがないときの周期 $\dfrac{2\pi}{4}=1.571$ 秒より、少しだけ長くなっています。ダンパーが動きにブレーキをかけるので、往復が少しおそくなるのです。$c=2$ では1.814秒（15\,\%長い）になります。くわしい式は6-12で求めます。

\kvihhead{対数減衰率}
となりの山の比の、逆数の自然対数を\textbf{対数減衰率}といい、$\delta$（デルタ）で表します。
\kvihframe{$\delta=\ln\dfrac{\text{前の山}}{\text{あとの山}}$}
$c=0.4$ では $\delta=\ln\dfrac{0.1}{0.0532}=\ln1.880\approx0.631$ です。$\delta$ が大きいほど、1回のゆれで大きくへります。\bkidx{たいすうげんすいりつ}{対数減衰率 $\delta$}

対数をとるのは、何回分かのへり方を「たし算」で考えられるからです。山が1回で0.532倍なら、$n$ 回で $0.532^n$ 倍。対数をとると $n\times0.631$ で、回数に比例します（3-12の $\log a^n=n\log a$）。また、$\delta=\alpha T_{\mathrm d}$（へる速さ×山と山の時間）なので、$\delta$ と $T_{\mathrm d}$ をはかれば、$\alpha$ がわかります。建物や橋の「ゆれの止まりやすさ」は、たたいてゆらし、山の高さをはかって、この $\delta$ から調べます。

\begin{seikaku}
山の高さが「同じ割合」でへるのは、ブレーキの力が速さに比例する（$-c\dot x$）からです。床の上をすべる物の摩擦のように、ブレーキの力の大きさが速さによらず一定のときは、山の高さは1回ごとに「同じ量」ずつへり（等差）、ある所でぴたりと止まります。やってみようの水の中や空気の中のゆれは、速さに比例する抵抗に近いので、割合でへりました。
\end{seikaku}

\bksubsection{5}{ちょうどよいブレーキをさがす：臨界と減衰比}

\label{bk:6-7:zeta}\bkhatenaref の問いにもどります。ブレーキ（ダンパー）は、どれくらいの強さがよいのでしょう。$m=0.5$、$k=8$ のまま、$c$ を少しずつかえて表計算し、2つのことを調べます。
\begin{itemize}[leftmargin=1.2\zw, itemsep=1pt, topsep=2pt, label={\color{\bkTT}●}]
\item 行きすぎ：0をこえて反対側へ、いちばん遠くまで行った距離（はじめの0.1に対する割合）
\item 落ち着く時刻：$\lvert x\rvert$ が0.002（はじめの2\,\%）より小さくなって、もう出なくなる時刻
\end{itemize}
\begin{bkcenter}
\begin{bktable}{colspec={Q[l,m,wd=30mm] *{8}{Q[c,m,wd=10.5mm]}}, rowsep=1pt, colsep=2pt}
$c$ & 0.4 & 1 & 2 & 3 & 3.5 & 4 & 5 & 10 \\
行きすぎ & 73\,\% & 44\,\% & 16\,\% & 2.8\,\% & 0.3\,\% & 0 & 0 & 0 \\
落ち着く時刻（s） & 9.59 & 3.53 & 2.02 & 1.43 & 1.11 & 1.46 & 2.10 & 4.74 \\
\end{bktable}
\bktabcaption{$c$ が4より小さいと行きすぎ、4以上では行きすぎない。}
\end{bkcenter}
行きすぎがなくなる境目は、$c=4$ です。$c=4$ のとき、行きすぎずに落ち着くまでの時間がいちばん短く（1.46秒）、それより強くすると、かえって時間がかかります（$c=10$ で4.74秒）。この境目の $c$ を\textbf{臨界減衰係数}といい、$c_{\mathrm c}$ と書きます。

\kvihhead{境目の $c$ のきまりを見つける}
$m$ と $k$ をかえて、同じように境目をさがします。
\begin{bkcenter}
\begin{bktable}{colspec={Q[c,m,wd=18mm] Q[c,m,wd=18mm] Q[c,m,wd=26mm] Q[c,m,wd=22mm] Q[c,m,wd=24mm]}, rowsep=1pt}
$m$ & $k$ & 境目の $c$ & $\sqrt{mk}$ & $2\sqrt{mk}$ \\
0.5 & 8 & 4 & 2 & 4 \\
0.5 & 2 & 2 & 1 & 2 \\
2 & 8 & 8 & 4 & 8 \\
1 & 25 & 10 & 5 & 10 \\
\end{bktable}
\bktabcaption{どの組でも、境目の $c$ は $2\sqrt{mk}$ に等しい。}
\end{bkcenter}
\kvihframe{$c_{\mathrm c}=2\sqrt{mk}$}
重い台車ほど、かたいばねほど、強いブレーキが必要になります。なぜ $2\sqrt{mk}$ になるのかは、6-12で、2次方程式の判別式から説明します（3-10で約束した「ゆれるか、ゆれないか」）。

\kvihhead{減衰比 $\zeta$}
使っている $c$ が、境目 $c_{\mathrm c}$ の何倍かを表す数を、\textbf{減衰比}といい、$\zeta$（ゼータ）で表します。\bkidx{げんすいひ}{減衰比 $\zeta$}
\kvihframe{$\zeta=\dfrac{c}{c_{\mathrm c}}=\dfrac{c}{2\sqrt{mk}}$}
$\zeta$ は「〔N·s/m〕÷〔N·s/m〕」なので、単位のない数です。$\zeta$ だけで、止まり方が3つに分かれます。
\begin{bkcenter}
\begin{bktable}{colspec={Q[c,m,wd=22mm] Q[l,m,wd=42mm] Q[l,m,wd=50mm]}, rowsep=1pt}
$\zeta$ & 止まり方 & ②のグラフ \\
$\zeta<1$ & ゆれながら止まる（減衰振動）& $c=0.4$（$\zeta=0.1$） \\
$\zeta=1$ & ちょうど（臨界） & $c=4$（$\zeta=1$） \\
$\zeta>1$ & ゆれずに、ゆっくりもどる & $c=10$（$\zeta=2.5$） \\
\end{bktable}
\bktabcaption{大学の本では、3つを「不足減衰」「臨界減衰」「過減衰」とよぶ（6-12）。}
\end{bkcenter}
$\zeta<1$ のゆれを、\textbf{減衰振動}といいます。この章の題の「ゆれながら止まる」動きです。\bkidx{げんすいしんどう}{減衰振動}

\kvihbox{\bkhatenaref の答え}
ブレーキは、$\zeta$ が1に近くなる強さにします。$\zeta$ が1より小さすぎると、針は目もりのまわりでゆれ続けます。$\zeta$ が1より大きすぎると、針はなかなか目もりに着きません。$\zeta=1$ が、行きすぎずにいちばん早く止まる強さです。表を見ると、$\zeta$ を1より少しだけ小さく（$c=3.5$、$\zeta=0.875$）すると、ほんの少し（0.3\,\%）行きすぎるかわりに、もっと早く（1.11秒）落ち着きます。計器や車では、目的に合わせて、$\zeta$ を1のまわりから選びます。

アナログの電流計では、針の根もとのコイルを巻いたアルミのわくが、磁石のあいだで回ります。わくが動くと、わくの中に\textbf{うず電流}が流れ、動きをじゃまする力が生まれます。この力は速さに比例するので、ダンパーと同じはたらきをします。

\kvihhead{$\zeta$ と、へる速さ・対数減衰率}
④の $\alpha=\dfrac{c}{2m}$ を、$\zeta$ と $\omega_{\mathrm n}$ で書きなおします。$c=2\zeta\sqrt{mk}$ なので、
\begin{bkhenkei}
\alpha=\frac{c}{2m}=\frac{2\zeta\sqrt{mk}}{2m} & $c=2\zeta\sqrt{mk}$ を入れる \\
=\zeta\sqrt{\frac{mk}{m^2}}=\zeta\sqrt{\frac{k}{m}} & $m=\sqrt{m^2}$ として、$\sqrt{\phantom{m}}$ をまとめる \\
=\zeta\,\omega_{\mathrm n} & $\omega_{\mathrm n}=\sqrt{\dfrac{k}{m}}$（6-6）\\
\end{bkhenkei}
$\zeta$ が小さいときは、山と山の時間 $T_{\mathrm d}$ は $\dfrac{2\pi}{\omega_{\mathrm n}}$ とほぼ同じなので、
\[ \delta=\alpha T_{\mathrm d}\approx\zeta\omega_{\mathrm n}\times\frac{2\pi}{\omega_{\mathrm n}}=2\pi\zeta \]
\begin{bkcenter}
\begin{bktable}{colspec={Q[l,m,wd=36mm] *{4}{Q[c,m,wd=17mm]}}, rowsep=1pt}
$\zeta$ & 0.05 & 0.1 & 0.25 & 0.5 \\
表計算の $\delta$ & 0.315 & 0.631 & 1.62 & 3.63 \\
$2\pi\zeta$ & 0.314 & 0.628 & 1.57 & 3.14 \\
\end{bktable}
\bktabcaption{$\zeta$ が0.1くらいまでなら、$\delta\approx2\pi\zeta$ はよく合う。}
\end{bkcenter}
だから、ゆれの山を2つはかれば、$\zeta\approx\dfrac{\delta}{2\pi}$ で減衰比がわかります。

\bkfukidashi{sora}{⑤は1回で全部わからなくても大丈夫。$\zeta=\dfrac{c}{2\sqrt{mk}}$ を計算して、1とくらべられれば先へ進めるよ。}

\bksubsection{6}{エネルギーは、どこへ行く？}

\label{bk:6-7:energy}6-6では、ばねのエネルギーと運動エネルギーの合計 $E=\dfrac{1}{2}mv^2+\dfrac{1}{2}kx^2$ が一定でした。ダンパーがあると、どうなるでしょう。6-6と同じように微分します。
\begin{bkhenkei}
\dot E=v\,(m\ddot x+kx) & 6-6と同じ計算 \\
=v\times(-c\dot x) & この章の式から $m\ddot x+kx=-c\dot x$ \\
=-cv^2 & $\dot x=v$ \\
\end{bkhenkei}
$c$ も $v^2$ も0以上なので、$\dot E$ はいつも0以下です。エネルギーは、ふえることがなく、動いているあいだずっとへっていきます。$cv^2$ は、ダンパーの力 $cv$ と速さ $v$ のかけ算です。力×1秒に進む距離なので、ダンパーが1秒あたりにする仕事（5-12の電力と同じ「1秒あたりのエネルギー」）です。

\noindent\begin{minipage}[c]{0.52\linewidth}
\begin{tikzpicture}
\begin{axis}[bkplot, width=60mm, height=42mm, xmin=0, xmax=6.4, ymin=0, ymax=0.046, xtick={2,4,6}, ytick={0.02,0.04}, yticklabels={0.02,0.04}, xlabel={$t$（s）}, ylabel={$E$（J）}]
  \addplot[bkPurple, line width=1.5pt] coordinates {(0.0,0.04) (0.25,0.03561) (0.5,0.02562) (0.75,0.02139) (1.0,0.01956) (1.25,0.01433) (1.5,0.01146) (1.75,0.01067) (2.0,0.00803) (2.25,0.00618) (2.5,0.00578) (2.75,0.00451) (3.0,0.00335) (3.25,0.00312) (3.5,0.00252) (3.75,0.00183) (4.0,0.00167) (4.25,0.00141) (4.5,0.00101) (4.75,0.00089) (5.0,0.00078) (5.25,0.00056) (5.5,0.00048) (5.75,0.00043) (6.0,0.00031)};
\end{axis}
\end{tikzpicture}
\end{minipage}\hfill
\begin{minipage}[c]{0.45\linewidth}
$c=0.4$ のときの $E$（表計算で0.25秒ごと）。台車が速く走る所でエネルギーが大きくへり、はしの近く（速さが0に近い所）では、ほとんどへりません。1周期（1.58秒）で、0.04\,Jが0.011\,Jまでへりました。
\end{minipage}\par\vspace{3pt}
へったエネルギーは、消えたのではありません。ダンパーの中で油が小さな穴をこすりながら流れ、油が温まります。ゆれのエネルギーは、\textbf{熱}に変わったのです。車で長いでこぼこ道を走ると、ショックアブソーバーがさわれないほど熱くなることがあります。

\bksubsection{7}{電気の減衰振動：RLC回路}

\label{bk:6-7:rlc}6-6のLC回路に、抵抗 $R$ を入れます。抵抗は、電気のエネルギーを熱に変える部品です（電気ストーブの線と同じ）。

\noindent\begin{minipage}[c]{0.4\linewidth}
\centering
\begin{tikzpicture}[x=0.9mm, y=0.9mm]
  \draw[bkInk, line width=0.8pt] (0,26) -- (36,26) (0,0) -- (36,0);
  \kvihcap{0}{0}{26}
  \kvihres{36}{13}{26}
  \kvihcoil{36}{0}{13}
  \node[font=\sffamily\footnotesize, anchor=east] at (-5,15.2) {$+q$};
  \node[font=\sffamily\footnotesize, anchor=east] at (-5,10.4) {$-q$};
  \node[font=\sffamily\footnotesize, anchor=west] at (39,19.5) {$R$};
  \node[font=\sffamily\footnotesize, anchor=west] at (39.5,6.5) {$L$};
  \node[font=\sffamily\footnotesize, anchor=north] at (0,-2) {$C$};
  \draw[bk arrow, bkOrange!85!black, line width=1pt] (11,29) -- (25,29);
  \node[font=\sffamily\footnotesize, anchor=south, text=bkOrange!80!black] at (18,29.5) {電流 $i$};
\end{tikzpicture}
\end{minipage}\hfill
\begin{minipage}[c]{0.57\linewidth}
はじめにコンデンサーに電気をため（上の板が $+q$）、コイルと抵抗につなぎます。6-6と同じように、電流 $i$ は上の板から出てコイルのほうへ流れる向きを＋とし、$i=-\dfrac{\dd q}{\dd t}$ です。

コンデンサーの電圧 $\dfrac{q}{C}$ が、抵抗の電圧 $Ri$（オームの法則）と、コイルの電圧 $L\dfrac{\dd i}{\dd t}$ に分かれます（6-4のRL回路と同じ考え）。
\end{minipage}
\begin{bkhenkei}
\frac{q}{C}=L\frac{\dd i}{\dd t}+Ri & 電圧が2つに分かれる \\
\frac{q}{C}=-L\ddot q-R\dot q & $i=-\dot q$、$\dfrac{\dd i}{\dd t}=-\ddot q$ を入れる \\
L\ddot q+R\dot q+\frac{1}{C}\,q=0 & 右の2つを左へ移す \\
\end{bkhenkei}
\kvihframe{$L\ddot q+R\dot q+\dfrac{1}{C}\,q=0$}
ばね・質量・ダンパーの式 $m\ddot x+c\dot x+kx=0$ と同じ形です。6-6の対応（$m\to L$、$k\to\dfrac{1}{C}$、$x\to q$）に、新しく\textbf{ダンパー $c$ → 抵抗 $R$} が加わりました。どちらも、エネルギーを熱に変える部品です。\bkidx{あーるえるしーかいろ}{RLC回路}

\kvihhead{臨界の抵抗と減衰比}
おきかえれば、臨界の値も減衰比も、そのまま使えます。
\[ R_{\mathrm c}=2\sqrt{L\times\frac{1}{C}}=2\sqrt{\frac{L}{C}},\qquad \zeta=\frac{R}{R_{\mathrm c}}=\frac{R}{2}\sqrt{\frac{C}{L}} \]
$L=\qty{10}{mH}$、$C=\qty{1}{\micro F}$ のとき、$\dfrac{L}{C}=\dfrac{0.01}{0.000001}=\num{10000}$ なので、$R_{\mathrm c}=2\sqrt{\num{10000}}=\qty{200}{\ohm}$ です。単位は $\sqrt{\mathrm{H}/\mathrm{F}}=\sqrt{\dfrac{\mathrm{V\,s/A}}{\mathrm{A\,s/V}}}=\sqrt{\dfrac{\mathrm{V}^2}{\mathrm{A}^2}}=\mathrm{V/A}=\Omega$ で、抵抗の単位になっています。

コンデンサーを\qty{10}{V}にためてから、$R=\qty{20}{\ohm}$（$\zeta=0.1$）、$\qty{200}{\ohm}$（$\zeta=1$）、$\qty{1}{k\ohm}$（$\zeta=5$）でつないだときの、コンデンサーの電圧です。表計算では、$\Delta t$ を\qty{0.00001}{s}（$\qty{10}{\micro s}$）にして、ばねと同じ型で計算できます。
\begin{bkfigure}
\begin{tikzpicture}
\begin{axis}[bkplot, width=124mm, height=54mm, xmin=0, xmax=3.2, ymin=-8.5, ymax=12, xtick={0.5,1,1.5,2,2.5,3}, xticklabels={}, ytick={-5,5,10}, xlabel={$t$（ms）}, ylabel={電圧（V）}]
  \addplot[bkcurve, domain=0:3.1, samples=300] {10*exp(-x)*(cos(deg(9.94987*x))+0.100504*sin(deg(9.94987*x)))};
  \addplot[bkcurve2, domain=0:3.1, samples=120] {10*(1+10*x)*exp(-10*x)};
  \addplot[bkPurple, line width=1.7pt, domain=0:3.1, samples=160] {10.1031*exp(-1.01021*x)-0.1031*exp(-98.9898*x)};
  \kvihxtl{0.5}{-8.5}{0.5}\kvihxtl{1}{-8.5}{1}\kvihxtl{1.5}{-8.5}{1.5}\kvihxtl{2}{-8.5}{2}\kvihxtl{2.5}{-8.5}{2.5}\kvihxtl{3}{-8.5}{3}
  \node[font=\sffamily\footnotesize, text=bkBlue, anchor=west, inner sep=1pt] at (axis cs:1.75,10.4) {$\qty{20}{\ohm}$（$\zeta=0.1$）：ゆれる};
  \node[font=\sffamily\footnotesize, text=bkOrange!80!black, anchor=west, inner sep=1pt] at (axis cs:1.75,8.2) {$\qty{200}{\ohm}$（$\zeta=1$）：ちょうど};
  \node[font=\sffamily\footnotesize, text=bkPurple, anchor=west, inner sep=1pt] at (axis cs:1.75,6.0) {$\qty{1}{k\ohm}$（$\zeta=5$）：ゆっくり};
\end{axis}
\end{tikzpicture}
\bkfigcaption{RLC回路のコンデンサーの電圧。ばねの3本のグラフ（③）と同じ3つの止まり方。ゆれの周期は約0.63\,ms。}
\end{bkfigure}
電気の回路でも、抵抗が小さいとゆれながら止まり、ちょうど $R_{\mathrm c}$ でゆれずに早く落ち着き、大きすぎるとゆっくり放電します。ばねのゆれが1秒ほどかかるのに対して、電気のゆれは1000分の1秒ほどで終わります。



\begin{senmonsho}[大学の教科書ふうの書き方]
\begin{bkhonbun}
1自由度の粘性減衰系の運動方程式 $m\ddot x+c\dot x+kx=0$ を $m$ で除し、固有角振動数 $\omega_{\mathrm n}=\sqrt{k/m}$ と減衰比 $\zeta=c/c_{\mathrm c}$（$c_{\mathrm c}=2\sqrt{mk}$ は臨界減衰係数）を用いると
\[ \ddot x+2\zeta\omega_{\mathrm n}\dot x+\omega_{\mathrm n}^{\,2}x=0 \]
を得る。$0<\zeta<1$ のとき不足減衰であり、自由振動の振幅は $e^{-\zeta\omega_{\mathrm n}t}$ に比例して減衰する。対数減衰率は $\delta=2\pi\zeta/\sqrt{1-\zeta^2}$ である。
\end{bkhonbun}
\bkyaku{「1自由度」は位置が1つの数 $x$ で決まること。「粘性減衰」は速さに比例するブレーキ。「$m$ で除し」は $m$ でわって。$\dfrac{c}{m}=2\zeta\omega_{\mathrm n}$、$\dfrac{k}{m}=\omega_{\mathrm n}^{\,2}$ と書きかえた形で、⑤の $\alpha=\zeta\omega_{\mathrm n}$ と同じ。「不足減衰」は $\zeta<1$ の「ゆれながら止まる」。$\delta$ のくわしい式は6-12で導く。}
\end{senmonsho}

\begin{motto}[4]
\kvihbox{なぜ $\alpha=\dfrac{c}{2m}$ なのか：エネルギーで考える}
ダンパーが弱いとき、1回のゆれのあいだは、ほとんど単振動（振幅 $X$）と同じに動きます。エネルギーは、はしでの値 $E=\dfrac{1}{2}kX^2$ です。

⑥より、エネルギーは1秒あたり $cv^2$ ずつへります。単振動の速度は $v=-X\omega_{\mathrm n}\sin(\omega_{\mathrm n}t)$ の形なので、$v^2$ の1周期の平均は、$\sin^2$ の平均が $\dfrac{1}{2}$（5-15）だから $\dfrac{1}{2}X^2\omega_{\mathrm n}^{\,2}$ です。$m\omega_{\mathrm n}^{\,2}=k$ を使うと、
\[ (cv^2\ \text{の平均})=\frac{c}{2}X^2\omega_{\mathrm n}^{\,2}=\frac{c}{m}\times\frac{1}{2}m\omega_{\mathrm n}^{\,2}X^2=\frac{c}{m}\times\frac{1}{2}kX^2=\frac{c}{m}E \]
つまり $\dot E\approx-\dfrac{c}{m}E$ で、エネルギーは「自分に比例してへる」（6-4）ので、$E\propto e^{-\frac{c}{m}t}$ です（$\propto$ は「比例する」）。エネルギーは振幅の2乗に比例するので、振幅は $\sqrt{E}$ に比例し、
\[ X\propto\sqrt{e^{-\frac{c}{m}t}}=e^{-\frac{c}{2m}t} \]
$\sqrt{\phantom{m}}$ をとると肩が半分になる（$\sqrt{e^{a}}=e^{\frac{a}{2}}$）ことが、2でわる理由でした。
\end{motto}

\begin{taisetsu}
\def\tsl#1{\makebox[5.4\zw][l]{\sffamily\bfseries\small\color{\bkTT}#1}}
\tsl{式}$m\ddot x+c\dot x+kx=0$\qquad（ダンパーの力は $-c\dot x$、速度と反対向き）\\[3pt]
\tsl{山のへり方}となりの山の比は一定。つつみ $e^{-\alpha t}$、$\alpha=\dfrac{c}{2m}=\zeta\omega_{\mathrm n}$\\[3pt]
\tsl{対数減衰率}$\delta=\ln\dfrac{\text{前の山}}{\text{あとの山}}\approx2\pi\zeta$（$\zeta$ が小さいとき）\\[3pt]
\tsl{減衰比}$\zeta=\dfrac{c}{2\sqrt{mk}}$\qquad $\zeta<1$ ゆれる、$\zeta=1$ ちょうど、$\zeta>1$ ゆっくり\\[3pt]
\tsl{RLC回路}$L\ddot q+R\dot q+\dfrac{1}{C}q=0$\qquad $\zeta=\dfrac{R}{2}\sqrt{\dfrac{C}{L}}$\\[3pt]
\tsl{エネルギー}ダンパーや抵抗が、1秒あたり $cv^2$ ずつ熱に変える。
\end{taisetsu}

\begin{kotoba}
\small
\bkword{ダンパー}{速さに比例して、動きと反対向きの力を出す部品。ショックアブソーバー、ドアクローザー。}
\bkword{減衰係数}{ダンパーの力 $-c\dot x$ の $c$〔N·s/m〕。}
\bkword{減衰振動}{ゆれながら、山がだんだん低くなって止まる動き。}
\bkword{対数減衰率}{となりの山の比の逆数の自然対数 $\delta$。}
\bkword{臨界・減衰比}{行きすぎがなくなる境目の $c_{\mathrm c}=2\sqrt{mk}$。$\zeta=\dfrac{c}{c_{\mathrm c}}$。}
\bkword{つつみ}{ゆれの山を結んだ曲線。減衰振動では $e^{-\alpha t}$ の形。}
\end{kotoba}

\bkkigocards{\bkkigocard{$\zeta$}{ゼータ}{減衰比。$\dfrac{c}{2\sqrt{mk}}$。1より小さいとゆれ、1でちょうど、1より大きいとゆっくりもどる}{damping ratio}{6-7}}
{\bkkigocard{$c$}{シー}{ダンパーの減衰係数〔N·s/m〕。6-5の2乗の抵抗の係数 $c$〔N·s$^2$/m$^2$〕とは別の量}{damping coefficient}{6-7}}

\begin{column}{ドアクローザーのひみつ}
学校や会社の重いドアの上についている箱が、ドアクローザーです。中には、ドアを閉める向きに引くばねと、油の入ったダンパーが入っています。ばねだけだと、ドアは「バタン！」と大きな音を立てて閉まり、指をはさむと危険です。ダンパーが、閉まる速さに比例したブレーキをかけるので、ドアはゆっくり、静かに閉まります。多くのドアクローザーには、油の通り道の太さを変えるねじがついていて、ねじを回すと $c$ が変わり、閉まる速さを調節できます。
\end{column}

%% ---------------- 例題 ----------------
\bksection{reidai}{例題でたしかめよう}

\begin{reidai}[式を立てて、加速度を求める]
車の1つのタイヤの上に、車体の\qty{300}{kg}がのっています。ばね定数は\qty{30000}{N/m}、ダンパーの減衰係数は\qty{1800}{N.s/m}です。車体が、つり合いの位置から上へ $x=\qty{0.02}{m}$ の所を、下向きに\qty{0.5}{m/s}で動いているとき（上向きを＋）、車体の加速度を求めましょう。\par
\kangaekata $m\ddot x=-kx-c\dot x$（②）。上向きが＋なので、$x=0.02$、$\dot x=-0.5$。ばねの力とダンパーの力を、符号に気をつけて1つずつ計算する。\par
\shiki
\kviha ばねの力：$-kx=-\num{30000}\times0.02=-600$（下向き）
\kviha ダンパーの力：$-c\dot x=-1800\times(-0.5)=+900$（上向き）
\kviha 合計：$-600+900=300$
\kviha 加速度：$\ddot x=\dfrac{300}{300}=1$\par
\kotae 上向きに\qty{1}{m/s^2}
\end{reidai}
下へ動いているので、ダンパーは上向きにじゃまをします。この瞬間は、ダンパーの力のほうがばねの力より大きいので、加速度は上向きです。

\begin{jibun}
$m=\qty{0.5}{kg}$、$k=\qty{8}{N/m}$、$c=\qty{0.4}{N.s/m}$ の台車が、$x=\qty{0.05}{m}$ の所を、右向きに\qty{0.2}{m/s}で動いています（右向きを＋）。加速度を求めましょう。
\bkjans{ばね $-8\times0.05=-0.4$、ダンパー $-0.4\times0.2=-0.08$。合計 $-0.48$。$\ddot x=\dfrac{-0.48}{0.5}=\qty{-0.96}{m/s^2}$（左向きに0.96）}
\end{jibun}

\begin{reidai}[表計算の2行を手で計算する]
$m=0.5$、$k=8$、$c=0.4$ の台車を、$x=0.1$、$v=0$ からはなします。手で計算しやすいように $\Delta t=0.05$ として、速度を先に進めるオイラー法で2歩進め、0.1秒後の位置と速度を求めましょう。\par
\kangaekata ③の順。①加速度 $a=\dfrac{-8x-0.4v}{0.5}$、②新しい速度 $v+a\Delta t$、③新しい位置 $x+(\text{新しい速度})\Delta t$。\par
\shiki
\kviha 1歩め：$a=\dfrac{-8\times0.1-0}{0.5}=-1.6$
\kvihb $v=0+(-1.6)\times0.05=-0.08$
\kvihb $x=0.1+(-0.08)\times0.05=0.096$
\kviha 2歩め：$a=\dfrac{-8\times0.096-0.4\times(-0.08)}{0.5}$
\kvihb $\phantom{a}=\dfrac{-0.768+0.032}{0.5}=-1.472$
\kvihb $v=-0.08+(-1.472)\times0.05=-0.1536$
\kvihb $x=0.096+(-0.1536)\times0.05=0.08832$\par
\kotae 0.1秒後に $x\approx\qty{0.0883}{m}$、$v\approx\qty{-0.154}{m/s}$
\end{reidai}

\begin{jibun}
$m=0.5$、$k=8$、$c=4$ の台車を、$x=0.1$、$v=0$ からはなします。$\Delta t=0.05$ で2歩進め、0.1秒後の位置と速度を求めましょう。
\bkjans{1歩め $a=-1.6$、$v=-0.08$、$x=0.096$。2歩め $a=\dfrac{-0.768+0.32}{0.5}=-0.896$、$v=-0.1248$、$x=0.08976$}
\end{jibun}

\begin{reidai}[となりの山の比から、対数減衰率]
減衰振動の山の高さが、\qty{0.1}{m}、\qty{0.0532}{m}、…と、同じ割合でへっていきます。対数減衰率 $\delta$ を求めましょう。また、山の高さがはじめの $\dfrac{1}{10}$ より小さくなるのは、何回めの山からですか（$\ln1.88=0.631$、$\ln10=2.303$）。\par
\kangaekata $\delta=\ln\dfrac{\text{前の山}}{\text{あとの山}}$（④）。$n$ 周期あとの山は $0.532^n$ 倍。対数をとると、$n$ 周期で $n\delta$ だけへる。\par
\shiki
\kviha $\dfrac{0.1}{0.0532}\approx1.88$ なので、$\delta=\ln1.88=0.631$
\kviha $n$ 周期あとの山がはじめの $\dfrac{1}{10}$ になるのは、$n\delta=\ln10$ のとき
\kviha $n=\dfrac{2.303}{0.631}\approx3.65$
\kviha $n$ は整数なので、4周期あと（$0.532^4\approx0.080$）。1回めの山から数えて5回め\par
\kotae $\delta\approx0.631$。5回めの山（4周期あと）から
\end{reidai}

\begin{jibun}
山の高さが\qty{20}{mm}、\qty{16}{mm}、…と、同じ割合でへるゆれがあります。対数減衰率を求めましょう。また、山の高さがはじめの $\dfrac{1}{10}$ より小さくなるのは、何周期あとですか（$\ln1.25=0.223$、$\ln10=2.303$）。
\bkjans{$\delta=\ln\dfrac{20}{16}=\ln1.25=0.223$。$n=\dfrac{2.303}{0.223}\approx10.3$ なので、11周期あと（$0.8^{11}\approx0.086$）}
\end{jibun}

\begin{reidai}[減衰比を計算して、止まり方を答える]
$m=\qty{0.5}{kg}$、$k=\qty{8}{N/m}$ の台車に、$c=0.4$、4、10 のダンパーをつけます。それぞれの減衰比 $\zeta$ を求め、止まり方を答えましょう。また、例題1の車（$m=300$、$k=\num{30000}$、$c=1800$）の $\zeta$ も求めましょう。\par
\kangaekata $c_{\mathrm c}=2\sqrt{mk}$、$\zeta=\dfrac{c}{c_{\mathrm c}}$（⑤）。1より小さいか、等しいか、大きいか。\par
\shiki
\kviha $c_{\mathrm c}=2\sqrt{0.5\times8}=2\sqrt4=4$
\kviha $\zeta=\dfrac{0.4}{4}=0.1$、\quad$\dfrac{4}{4}=1$、\quad$\dfrac{10}{4}=2.5$
\kviha 車：$c_{\mathrm c}=2\sqrt{300\times30000}=2\sqrt{9000000}=2\times3000=6000$
\kvihb $\zeta=\dfrac{1800}{6000}=0.3$\par
\kotae 0.1（ゆれながら止まる）、1（ちょうど）、2.5（ゆっくりもどる）。車は0.3（少しゆれながら止まる）
\end{reidai}

\begin{jibun}
$m=\qty{2}{kg}$、$k=\qty{200}{N/m}$ のおもりに、$c=10$、40、60 のダンパーをつけます。それぞれの $\zeta$ を求め、止まり方を答えましょう。
\bkjans{$c_{\mathrm c}=2\sqrt{400}=40$。$\zeta=0.25$（ゆれる）、1（ちょうど）、1.5（ゆっくり）}
\end{jibun}

\begin{reidai}[ほしい減衰比にするダンパーを選ぶ]
例題1の車（$m=\qty{300}{kg}$、$k=\qty{30000}{N/m}$）で、減衰比を $\zeta=0.4$ にしたいと思います。ダンパーの減衰係数 $c$ をいくつにすればよいですか。また、このときの固有角振動数 $\omega_{\mathrm n}$ と、へる速さ $\alpha$ を求めましょう。\par
\kangaekata $\zeta=\dfrac{c}{c_{\mathrm c}}$ を $c$ について解くと $c=\zeta\times c_{\mathrm c}$。$\alpha=\zeta\omega_{\mathrm n}$（⑤）。\par
\shiki
\kviha $c_{\mathrm c}=6000$（例題4）なので、$c=0.4\times6000=2400$
\kviha $\omega_{\mathrm n}=\sqrt{\dfrac{30000}{300}}=\sqrt{100}=10$
\kviha $\alpha=0.4\times10=4$（確かめ：$\dfrac{c}{2m}=\dfrac{2400}{600}=4$）\par
\kotae $c=\qty{2400}{N.s/m}$、$\omega_{\mathrm n}=10$\,rad/s、$\alpha=4$\,/s
\end{reidai}

\begin{jibun}
$m=\qty{400}{kg}$、$k=\qty{40000}{N/m}$ の車で、$\zeta=0.25$ にしたいとき、$c$ と $\omega_{\mathrm n}$、$\alpha$ を求めましょう。
\bkjans{$c_{\mathrm c}=2\sqrt{16000000}=8000$。$c=0.25\times8000=\qty{2000}{N.s/m}$。$\omega_{\mathrm n}=\sqrt{100}=10$\,rad/s、$\alpha=2.5$\,/s}
\end{jibun}
同じ車に荷物をたくさん積むと、$m$ がふえます。$c_{\mathrm c}=2\sqrt{mk}$ が大きくなるので、同じダンパーでも $\zeta=\dfrac{c}{c_{\mathrm c}}$ は小さくなり、車はゆれやすくなります。トラックやバスには、積む荷物の重さに合わせて、ばねやダンパーの強さを変えられるものもあります。

\begin{reidai}[つつみから、ゆれがおさまる時間を求める]
例題5の車（$\alpha=4$\,/s）が、でこぼこを通って\qty{5}{cm}ゆれ始めました。ゆれのつつみが\qty{5}{mm}（はじめの $\dfrac{1}{10}$）になるのは何秒後ですか。また、$\alpha$ が半分の2\,/sなら何秒後ですか（$\ln10=2.303$）。\par
\kangaekata つつみは $X_0e^{-\alpha t}$（④）。$e^{-\alpha t}=\dfrac{1}{10}$ になる $t$ を、両辺の自然対数をとって求める（5-14）。\par
\shiki
\kviha $e^{-4t}=\dfrac{1}{10}$
\kviha $-4t=\ln\dfrac{1}{10}=-\ln10=-2.303$（両辺の自然対数をとる）
\kviha $t=\dfrac{2.303}{4}\approx0.58$
\kviha $\alpha=2$ なら $t=\dfrac{2.303}{2}\approx1.15$\par
\kotae 約\qty{0.58}{s}後。$\alpha=2$ なら約\qty{1.15}{s}後（2倍かかる）
\end{reidai}
$\ln\dfrac{1}{10}=-\ln10$ は、5-18の冷めるお湯の $\ln\dfrac{1}{2}=-\ln2$ と同じきまりです。へる速さ $\alpha$ が半分になると、おさまるまでの時間は2倍になります。

\begin{jibun}
へる速さ $\alpha=0.4$\,/s の台車（③の $c=0.4$）のゆれのつつみが、はじめの $\dfrac{1}{10}$ になるのは何秒後ですか。また、$\dfrac{1}{100}$ になるのは何秒後ですか（$\ln10=2.303$、$\ln100=4.605$）。
\bkjans{$t=\dfrac{2.303}{0.4}\approx\qty{5.8}{s}$後。$\dfrac{1}{100}$ は $\dfrac{4.605}{0.4}\approx\qty{11.5}{s}$後}
\end{jibun}

\begin{reidai}[山の比から減衰比を見積もる]
高いビルを、強い風のあとで計ってみると、ゆれの山が1回ごとに0.9倍になっていました。対数減衰率と、減衰比 $\zeta$ を見積もりましょう（$\ln\dfrac{1}{0.9}=0.105$、$\pi=3.14$）。\par
\kangaekata $\delta=\ln\dfrac{\text{前の山}}{\text{あとの山}}=\ln\dfrac{1}{0.9}$。$\zeta$ が小さいときは $\zeta\approx\dfrac{\delta}{2\pi}$（⑤）。\par
\shiki
\kviha $\delta=\ln\dfrac{1}{0.9}=0.105$
\kviha $\zeta\approx\dfrac{0.105}{2\times3.14}=\dfrac{0.105}{6.28}\approx0.017$\par
\kotae $\delta\approx0.105$、$\zeta\approx0.017$（約2\,\%）
\end{reidai}
ビルの $\zeta$ は、ふつう0.01〜0.05ほどととても小さく、一度ゆれ始めると、なかなか止まりません。そこで、超高層ビルでは、柱やかべの中にダンパーを入れて、$\zeta$ を大きくします（物理・工学ではこう使う）。

\begin{jibun}
橋をたたいてゆらすと、ゆれの山が1回ごとに0.8倍になりました。対数減衰率と減衰比を見積もりましょう（$\ln1.25=0.223$、$\pi=3.14$）。
\bkjans{$\delta=\ln\dfrac{1}{0.8}=\ln1.25=0.223$。$\zeta\approx\dfrac{0.223}{6.28}\approx0.036$}
\end{jibun}

\begin{reidai}[RLC回路の臨界の抵抗と減衰比]
$L=\qty{10}{mH}$、$C=\qty{1}{\micro F}$ のRLC回路の、臨界の抵抗 $R_{\mathrm c}$ を求めましょう。また、$R=\qty{20}{\ohm}$ のときの減衰比を求め、止まり方を答えましょう。\par
\kangaekata $R_{\mathrm c}=2\sqrt{\dfrac{L}{C}}$、$\zeta=\dfrac{R}{R_{\mathrm c}}$（⑦）。接頭語を直してから計算する。\par
\shiki
\kviha $L=0.01$、$C=0.000001$ なので、$\dfrac{L}{C}=\dfrac{0.01}{0.000001}=\num{10000}$
\kviha $R_{\mathrm c}=2\sqrt{10000}=2\times100=200$
\kviha $\zeta=\dfrac{20}{200}=0.1$\par
\kotae $R_{\mathrm c}=\qty{200}{\ohm}$。$\zeta=0.1$ で、電圧はゆれながら0になる
\end{reidai}

\begin{jibun}
$L=\qty{1}{mH}$、$C=\qty{10}{\micro F}$ のRLC回路の、臨界の抵抗を求めましょう。また、$R=\qty{5}{\ohm}$ のときの減衰比を求め、止まり方を答えましょう。
\bkjans{$\dfrac{L}{C}=\dfrac{0.001}{0.00001}=100$。$R_{\mathrm c}=2\times10=\qty{20}{\ohm}$。$\zeta=\dfrac{5}{20}=0.25$ で、ゆれながら止まる}
\end{jibun}
\noindent\begin{minipage}[c]{0.5\linewidth}
\centering
\begin{tikzpicture}[x=1mm, y=1mm]
  \fill[bkBlue!12] (0,-2) rectangle (24,2);
  \fill[bkPurple!12] (24,-2) rectangle (60,2);
  \draw[bkInk, line width=0.7pt, -{Stealth[length=4pt]}] (0,0) -- (62,0) node[anchor=west, font=\sffamily\footnotesize] {$\zeta$};
  \foreach \u/\l in {0/0, 24/1, 48/2} {\draw[bkInk, line width=0.6pt] (\u,-1.2) -- (\u,1.2); \node[font=\sffamily\footnotesize, anchor=north] at (\u,-2.6) {\l};}
  \fill[bkBlue] (2.4,0) circle (1.2);
  \fill[bkOrange] (24,0) circle (1.2);
  \node[font=\sffamily\footnotesize, anchor=south, text=bkBlue] at (4.5,3) {$\qty{20}{\ohm}$};
  \node[font=\sffamily\footnotesize, anchor=south, text=bkOrange!80!black] at (24,3) {$\qty{200}{\ohm}$};
  \node[font=\sffamily\footnotesize, anchor=north, text=bkInk] at (12,-7) {ゆれる};
  \node[font=\sffamily\footnotesize, anchor=north, text=bkInk] at (42,-7) {ゆっくり};
\end{tikzpicture}
\end{minipage}\hfill
\begin{minipage}[c]{0.47\linewidth}
例題8の回路で、抵抗を $\zeta$ の数直線にのせました。$\zeta=1$（$\qty{200}{\ohm}$）が境目です。$\qty{1}{k\ohm}$ は $\zeta=5$ で、ずっと右の「ゆっくり」の側にあります。
\end{minipage}

\begin{reidai}[ダンパーが熱に変えたエネルギー]
$k=\qty{8}{N/m}$ のばねの台車が、山の高さ\qty{0.1}{m}から、1周期あとの山の高さ\qty{0.0532}{m}までゆれました。このあいだにダンパーが熱に変えたエネルギーは何Jですか。それは、はじめのエネルギーの何\,\%ですか。\par
\kangaekata 山（はし）では速さが0なので、エネルギーはばねの $\dfrac{1}{2}kx^2$ だけ（⑥）。へった分が、熱に変わった分。\par
\shiki
\kviha はじめ：$\dfrac{1}{2}\times8\times0.1^2=0.04$
\kviha 1周期あと：$\dfrac{1}{2}\times8\times0.0532^2=4\times0.00283=0.0113$
\kviha 熱に変わった分：$0.04-0.0113=0.0287$
\kviha 割合：$0.0287\div0.04\approx0.72$\par
\kotae 約\qty{0.029}{J}。はじめの約72\,\%
\end{reidai}
山の高さは0.532倍ですが、エネルギーは $0.532^2\approx0.28$ 倍です。エネルギーは高さの2乗に比例するからです。

\begin{jibun}
$k=\qty{50}{N/m}$ のばねのおもりの山の高さが、\qty{0.2}{m}から1周期あとに\qty{0.1}{m}になりました。熱に変わったエネルギーと、その割合を求めましょう。
\bkjans{はじめ $\frac{1}{2}\times50\times0.04=\qty{1}{J}$、あと $\frac{1}{2}\times50\times0.01=\qty{0.25}{J}$。熱に変わったのは\qty{0.75}{J}で75\,\%}
\end{jibun}

%% ---------------- まちがえやすいところ ----------------
\bksection{tsumazuki}{まちがえやすいところをチェック}

\begin{tsumazuki}{ダンパーの力を、位置と反対向きにする}
\bkng{右のはし（$x>0$）では、ダンパーが左へいちばん強く引く}
\bkok{右のはしでは止まっている（$\dot x=0$）ので、ダンパーの力は0}
ダンパーの力は $-c\dot x$ で、速度と反対向きです。位置と反対向きなのは、ばねの力 $-kx$ のほうです。
\end{tsumazuki}

\begin{tsumazuki}{ブレーキが強いほど、早く止まると思う}
\bkng{$c=10$（$\zeta=2.5$）のほうが、$c=4$ より早く止まる}
\bkok{いちばん早く落ち着くのは、$\zeta$ が1の近く}
強すぎるダンパーは、もどる動きもじゃまします。③のグラフで、$c=10$ は2秒たっても、まだ約\qty{2}{cm}残っています（$c=4$ ではもう\qty{0.3}{mm}）。
\end{tsumazuki}

\begin{tsumazuki}{臨界の $c$ で2をかけ忘れる}
\bkng{$c_{\mathrm c}=\sqrt{mk}=\sqrt{0.5\times8}=2$}
\bkok{$c_{\mathrm c}=2\sqrt{mk}=2\times2=4$}
⑤の表で、境目はどれも $\sqrt{mk}$ の2倍でした。$c=2$ のグラフはまだ16\,\%も行きすぎます。
\end{tsumazuki}

\begin{tsumazuki}{対数減衰率で、比を逆にする}
\bkng{$\delta=\ln\dfrac{0.0532}{0.1}=\ln0.532=-0.631$}
\bkok{$\delta=\ln\dfrac{0.1}{0.0532}=\ln1.88=0.631$}
$\delta$ は「前の山÷あとの山」（1より大きい）の対数で、正の数になるように決めてあります。逆にすると、符号が反対になります。
\end{tsumazuki}

\begin{tsumazuki}{6-5の $c$ と同じ量だと思う}
\bkng{ダンパーの力は $cv^2$}
\bkok{ダンパーの力は $c\dot x$（速さに比例）。単位は〔N·s/m〕}
6-5の $c$ は2乗の抵抗の係数〔N·s$^2$/m$^2$〕でした。同じ文字でも、単位のちがう別の量です。
\end{tsumazuki}



\begin{column}{計算のせいでふくらむ？}
6-2のふつうのオイラー法（今の速度で位置を進める）で $c=0.4$ のゆれを計算すると、$\Delta t=0.1$ では、山は0.1から0.19、0.34、0.64、…と、へるどころか1往復ごとに約1.9倍に高くなってしまいます。計算の方法がエネルギーを少しずつふやす分（6-2）が、ダンパーがへらす分より大きいからです。これを見て「ダンパーが効いていない」と思ったら、まちがいです。計算の方法を、速度を先に進めるオイラー法にかえるか、$\Delta t$ を小さくして確かめます。表計算の答えをうのみにせず、エネルギーがふえていないかを見るのは、技術者の大事な習慣です。
\end{column}

%% ---------------- 物理・工学 ----------------
\begin{butsuri}
\begin{itemize}[leftmargin=1.2\zw, itemsep=2pt, topsep=0pt, label={\color{bkPurple}●}]
\item \textbf{車のサスペンション}\par\noindent
\begin{minipage}[c]{0.64\linewidth}
タイヤと車体のあいだには、ばねとダンパー（ショックアブソーバー）がならんでいます。ばねがでこぼこの衝撃をやわらげ、ダンパーがそのあとのゆれを止めます。乗用車では、減衰比を0.3前後にすることが多く、少しだけゆれるかわりに、でこぼこの衝撃を体に伝えにくくしています。
\end{minipage}\hfill
\begin{minipage}[c]{0.32\linewidth}
\centering
\begin{tikzpicture}[x=0.9mm, y=0.9mm]
  \filldraw[fill=bkBlue!20, draw=bkBlue, line width=0.8pt, rounded corners=1mm] (-18,22) rectangle (18,30);
  \kvihspringv{-6}{22}{8}
  \kvihdamperv{6}{22}{8}
  \fill[bkInk!80] (-12,8) rectangle (12,6);
  \filldraw[fill=bkInk!70, draw=bkInk, line width=0.6pt] (0,0) circle (6);
  \fill[bkGray!60] (0,0) circle (2.2);
  \draw[bkGray, line width=1pt] (-20,-6) -- (20,-6);
  \node[font=\sffamily\footnotesize] at (0,26) {車体};
\end{tikzpicture}
\end{minipage}
\item \textbf{超高層ビルの制振ダンパー}\par\noindent
\begin{minipage}[c]{0.64\linewidth}
ビルの柱と柱のあいだに、ななめにオイルダンパーを入れます。地震や強い風でビルが横にゆれると、ダンパーがのび縮みして、ゆれのエネルギーを熱に変えます。$\zeta$ がわずか0.02ほどのビルでも、0.1ほどまで大きくでき、ゆれが早くおさまります。
\end{minipage}\hfill
\begin{minipage}[c]{0.32\linewidth}
\centering
\begin{tikzpicture}[x=1mm, y=1mm]
  \draw[bkInk, line width=1pt] (-12,0) -- (-12,30) (12,0) -- (12,30) (-12,15) -- (12,15) (-12,30) -- (12,30);
  \draw[bkGray, line width=1.2pt] (-16,0) -- (16,0);
  \draw[bkYellow!70!black, line width=2.6pt] (-11,1) -- (-1,7);
  \draw[bkInk!80, line width=1pt] (-1,7) -- (11,14);
  \draw[bkYellow!70!black, line width=2.6pt] (-11,16) -- (-1,22);
  \draw[bkInk!80, line width=1pt] (-1,22) -- (11,29);
  \node[font=\sffamily\footnotesize, anchor=north] at (0,-1.5) {ななめのダンパー};
\end{tikzpicture}
\end{minipage}
\item \textbf{電源を入れたときの電圧のゆれ}\par\noindent 電子回路の電源の配線には、コイルとコンデンサーのはたらきが自然にあり、抵抗が小さいと、スイッチを入れた直後に電圧がゆれます（リンギング）。半導体の部品は、決められた電圧をこえるとこわれることがあるので、抵抗を入れて $\zeta$ を大きくし、ゆれをおさえます。
\end{itemize}
\end{butsuri}

%% ---------------- 章末問題 ----------------
\bkendproblems
\bkproblemgroup{基本}{3}

\bkproblem $m=\qty{2}{kg}$、$k=\qty{50}{N/m}$、$c=\qty{4}{N.s/m}$ のおもりが、$x=\qty{0.1}{m}$ の所を、左向きに\qty{0.5}{m/s}で動いています（右向きを＋）。加速度を求めましょう。
\bkpans{ばね $-5$、ダンパー $-4\times(-0.5)=+2$。合計 $-3$。$\ddot x=\qty{-1.5}{m/s^2}$}

\bkproblem $m=0.5$、$k=8$、$c=2$ の台車を、$x=0.1$、$v=0$ からはなします。$\Delta t=0.05$ で、速度を先に進めるオイラー法で2歩進めましょう。
\bkpans{1歩め $a=-1.6$、$v=-0.08$、$x=0.096$。2歩め $a=\dfrac{-0.768+0.16}{0.5}=-1.216$、$v=-0.1408$、$x=0.08896$}

\bkproblem 山の高さが\qty{50}{mm}、\qty{40}{mm}、…と同じ割合でへるゆれの、対数減衰率を求めましょう。山がはじめの $\dfrac{1}{10}$ より小さくなるのは、何周期あとですか（$\ln1.25=0.223$、$\ln10=2.303$）。
\bkpans{$\delta=0.223$。$\dfrac{2.303}{0.223}\approx10.3$ なので11周期あと}

\bkproblem $m=\qty{1}{kg}$、$k=\qty{100}{N/m}$ のおもりの臨界減衰係数を求めましょう。$c=5$、20、30 のときの $\zeta$ と止まり方を答えましょう。
\bkpans{$c_{\mathrm c}=2\sqrt{100}=20$。$\zeta=0.25$（ゆれる）、1（ちょうど）、1.5（ゆっくり）}

\bkproblem $L=\qty{4}{mH}$、$C=\qty{1}{\micro F}$ のRLC回路の、臨界の抵抗を求めましょう。$R=\qty{40}{\ohm}$ のときの $\zeta$ も求めましょう。
\bkpans{$\dfrac{L}{C}=\dfrac{0.004}{0.000001}=4000$、$R_{\mathrm c}=2\sqrt{4000}\approx\qty{126}{\ohm}$。$\zeta=\dfrac{40}{126}\approx0.32$}

\bkproblem $k=\qty{20}{N/m}$ のばねのおもりの山の高さが、\qty{0.1}{m}から1周期あとに\qty{0.06}{m}になりました。熱に変わったエネルギーと、その割合を求めましょう。
\bkpans{はじめ\qty{0.1}{J}、あと $\frac{1}{2}\times20\times0.0036=\qty{0.036}{J}$。\qty{0.064}{J}で64\,\%}



\bkproblemgroup{チャレンジ}{4}

\bkproblem 車体\qty{250}{kg}、ばね定数\qty{25000}{N/m}の車で、減衰比を0.3にしたいとき、ダンパーの減衰係数と、へる速さ $\alpha$ を求めましょう。
\bkpans{$c_{\mathrm c}=2\sqrt{6250000}=5000$、$c=\qty{1500}{N.s/m}$。$\omega_{\mathrm n}=10$、$\alpha=3$\,/s}

\bkproblem 塔をゆらすと、ゆれの山が1回ごとに0.95倍になりました。減衰比を見積もりましょう。また、山がはじめの半分より小さくなるのは何周期あとですか（$\ln\dfrac{1}{0.95}=0.0513$、$\ln2=0.693$、$\pi=3.14$）。
\bkpans{$\delta=0.0513$、$\zeta\approx\dfrac{0.0513}{6.28}\approx0.008$。$\dfrac{0.693}{0.0513}\approx13.5$ なので14周期あと}

\bkproblem $L=\qty{10}{mH}$、$C=\qty{1}{\micro F}$ のRLC回路で、減衰比を0.5にするには、抵抗を何 $\Omega$ にすればよいですか。また、$R=\qty{300}{\ohm}$ にすると、止まり方はどうなりますか。
\bkpans{$R_{\mathrm c}=\qty{200}{\ohm}$ なので $R=0.5\times200=\qty{100}{\ohm}$。$\qty{300}{\ohm}$ では $\zeta=1.5$ で、ゆれずにゆっくり放電する}

\bkproblem ゆれのつつみ（$e^{-\alpha t}$ の形）が小さくなるまでの時間を求めましょう（$\ln10=2.303$、$\ln100=4.605$）。\newline（1）へる速さが $\alpha=1.5$\,/s のとき、つつみがはじめの $\dfrac{1}{10}$ になるのは何秒後ですか。$\dfrac{1}{100}$ になるのは何秒後ですか。\newline（2）$m=\qty{0.5}{kg}$、$k=\qty{8}{N/m}$ の台車に、$c=\qty{1}{N.s/m}$ のダンパーをつけました。減衰比 $\zeta$ と、へる速さ $\alpha$ を求め、つつみがはじめの $\dfrac{1}{10}$ になるのは何秒後か求めましょう。
\bkpans{（1）\nobreak $\dfrac{2.303}{1.5}\approx\qty{1.5}{s}$後、$\dfrac{4.605}{1.5}\approx\qty{3.1}{s}$後。（2）\nobreak $\zeta=\dfrac{1}{4}=0.25$。$\alpha=\dfrac{1}{2\times0.5}=1$\,/s。$\dfrac{2.303}{1}\approx\qty{2.3}{s}$後}

\begin{furikaeri}
\begin{checklist}
\item ばね・質量・ダンパーの式を、力の図から立てられる。\bkfurref{bk:6-7:shiki}
\item となりの山の比から、対数減衰率を求められる。\bkfurref{bk:6-7:yama}
\item 減衰比 $\zeta$ を計算して、止まり方を答えられる（RLC回路でも）。\bkfurref{bk:6-7:zeta}
\end{checklist}
\end{furikaeri}

\bkfukidashi{delta}{$\zeta=1$ がちょうどよいブレーキ。でも、なぜ $2\sqrt{mk}$ なの？\quad そのなぞは、6-12で2次方程式が教えてくれるよ。}
